\title{Large Language Models Can Automatically Engineer Features for Few-Shot Tabular Learning}

\begin{abstract}
Large Language Models (LLMs) have demonstrated remarkable capabilities in addressing complex and previously unencountered reasoning challenges, rendering them exceptionally promising for the field of tabular learning—a domain critical to various practical settings. This work introduces an innovative in-context learning methodology called \textsf{FeatLLM}, where LLMs operate as feature constructors to devise input datasets precisely tailored for superior performance in tabular prediction tasks. The features created via \textsf{FeatLLM} enable downstream models, such as linear regression, to effectively estimate class probabilities, yielding strong results in few-shot scenarios. Importantly, at inference, our framework depends solely on this straightforward predictive model utilizing the derived features, bypassing the LLM entirely. Distinct from existing approaches leveraging LLMs, \textsf{FeatLLM} obviates the necessity of querying the LLM during inference on each individual instance. The system also functions with just API-level LLM access, inherently sidestepping limitations imposed by prompt length. Empirical results on a wide array of tabular datasets reveal that \textsf{FeatLLM} generates robust and interpretable rules, offering a notable ($10\%$ mean) improvement over methods like TabLLM and STUNT.
\end{abstract}

\section{Introduction}
In recent developments, Large Language Models (LLMs) have displayed significant proficiency in producing appropriate solutions for novel, previously untrained tasks, all without requiring fine-tuning specific to those tasks~\cite{creswell2022selection,imani2023mathprompter,taylor2022galactica}. Owing to pretraining on large-scale textual data, LLMs accumulate extensive background knowledge, which can substitute manual domain expertise in a variety of contexts~\cite{Genomes2021,singhal2023large,wilcox2003role}. Their deployment is increasingly central to automating tasks in diverse fields, including, but not limited to, dialogue comprehension~\cite{finch2023leveraging} and code base debugging~\cite{fan2023automated}. Notably, when equipped with a brief description of a task and several relevant examples in the prompt, LLMs are capable of grasping contextual information and placing emphasis on crucial features or instances~\cite{brown2020language,zhao2023survey}, evidencing their facility for data-driven reasoning and their utility as advanced feature generators.

Leveraging the inherent reasoning strengths and accumulated knowledge within LLMs has extended into tabular data analysis. For example, insights such as the higher susceptibility of older populations to specific diseases can directly inform predictions within disease classification tasks. In these scenarios, cutting-edge methods encode both task and variable information in natural language, serialize tabular inputs, and provide this structured text to the LLM~\cite{dinh2022lift,wang2023anypredict}. Subsequently, the workflow often involves techniques like in-context learning~\cite{nam2023semi} or the application of parameter-efficient adaptation strategies~\cite{dinh2022lift,hegselmann2023tabllm}. In empirical studies, harnessing LLM-based prior knowledge has yielded substantial gains over traditional tabular learning algorithms in settings where data is scarce~\cite{hegselmann2023tabllm}.

Existing LLM-driven methods for learning from tabular data are confronted with several notable shortcomings. In direct end-to-end inference scenarios, each instance necessitates an individual LLM evaluation, thereby incurring substantial computational overhead. Additionally, achieving competitive accuracy generally requires fine-tuning the LLM, an unrealistic demand when using models that lack publicly accessible training infrastructure or tuning interfaces—a relevant concern as recent state-of-the-art LLMs predominantly restrict users to inference-only APIs~\cite{anil2023palm,openai2023gpt4}. Moreover, most available methods struggle with extended prompt lengths~\cite{liu2023lost}, particularly when the tabular dataset includes numerous features such that the prompt length swells, which commonly renders the approach impractical or reduces prediction quality. Collectively, these limitations inhibit the applicability of LLM-based tabular learning approaches in practical settings.

To address these issues, our method reframes the involvement of LLMs, utilizing them primarily for feature construction rather than direct prediction. In this paradigm, the LLM leverages its rich background knowledge to enhance feature selection. Concretely, we introduce an innovative framework, \textsf{FeatLLM}, which employs in-context learning to uncover the underlying “criteria” used by the LLM to make predictions. For example, when classifying disease status, the LLM can interpret and output a set of rules—criteria based on feature values—that dictate disease identification. These synthesized criteria are then used to construct new features, offering a substitution for the original features during downstream tasks. This repositioning enables the LLM to focus on feature engineering, simplifying subsequent model requirements; once the refined features are established, downstream predictions can be executed by straightforward machine learning algorithms, thereby minimizing inference time. Harnessing the in-context learning capabilities of LLMs, features are discoverable without requiring model retraining—demonstrations can simply be incorporated into prompt examples. In addition, techniques such as feature bagging, inspired by ensemble learning methods~\cite{breiman2001random}, help manage prompt size constraints, facilitating scalability.

\textsf{FeatLLM} formulates the input prompt for feature generation into two distinct phases: (1) interpreting the task and identifying how features correlate with target variables; and (2) formulating definitive rules to separate classes, grounded in insights from the initial phase and supplemented by a handful of examples. This sequential reasoning guides large language models (LLMs) in dismissing irrelevant features during rule derivation. The inferred rules are then implemented on the dataset (with rules processed as programmatic instructions), resulting in a transformation that encodes each instance as a binary feature reflecting whether the rule is satisfied. Subsequently, a simple model is trained to predict class probabilities using these synthesized features. The framework's reliability is enhanced through an ensemble mechanism that repeatedly generates features along with feature bagging. Careful tuning of the selected feature and sample quantities during bagging addresses the issue of oversized prompts. \textsf{FeatLLM} operates independently of training requirements and incurs minimal inference costs, making LLM integration practical across diverse applications.

We assess \textsf{FeatLLM} on 13 separate tabular datasets within low-shot scenarios, demonstrating robust and competitive performance. Our results indicate that LLMs effectively leverage both existing domain knowledge and information extracted from the data, producing valuable features. Across multiple configurations, our method consistently surpasses leading few-shot learning approaches. The implementation is made available through an anonymized GitHub repository at~\texttt{https://github.com/Sungwon-Han/FeatLLM}.

\section{Related Work}

\subsection{Few-Shot Learning with Tabular Data} Progress in few-shot learning has made it possible to derive broadly applicable data representations with only a handful of labeled examples~\cite{chen2018closer}. While initial developments in this area centered on image data, subsequent research has established the generalizability of few-shot approaches across multiple modalities such as text, audio, and more recently, tabular datasets~\cite{majumder2022few,nam2023stunt,schick2021exploiting}. Relying on limited labeled data heightens the possibility of picking up on incidental correlations~\cite{bartlett2021deep}. Accordingly, modern strategies integrate supplementary information sources or embed prior knowledge specific to the domain to introduce effective inductive biases into model training. For instance, several works have explored leveraging abundant unlabeled samples through thoughtful data augmentation under semi-supervised paradigms~\cite{ucar2021subtab,yoon2020vime}, or utilizing unsupervised pre-training methods that are not tied to any specific downstream task~\cite{somepalli2021saint}. Common unsupervised training techniques may mask or intentionally corrupt input portions for subsequent recovery~\cite{arik2021tabnet,majmundar2022met}, or utilize augmented views of data to facilitate contrastive objectives~\cite{bahri2022scarf,chen2020simple}. Nonetheless, the diversity inherent in tabular data often hinders the effectiveness of generic augmentation strategies, limiting their success in few-shot scenarios~\cite{nam2023stunt}. 

Emerging methodologies advocate training on an extensive variety of tabular datasets, extending well beyond direct relevance to a given task, and subsequently adapting to the intended application via transfer learning~\cite{zhang2023generative,levin2022transfer,zhu2023xtab}. As an illustration, TransTab~\cite{wang2022transtab} utilizes transformers that acquire knowledge of inter-column semantics not only by interpreting table entries but also through analysis of column headings. This approach allows the system to leverage the accumulated semantic insights for zero-shot or few-shot learning tasks on novel tabular data. Alternatively, TabPFN~\cite{hollmann2022tabpfn} constructs pre-training data by blending distributions of various types, pre-training a Bayesian neural network using this synthetic data, and subsequently performing inference by conditioning on both train and test points for the target dataset, with no additional fine-tuning required. The generalized knowledge amassed across these diverse sources has enabled TabPFN to outperform conventional tree-ensemble methods and demonstrated efficacy in low-data regimes.

\subsection{Language-Interfaced Tabular Learning} An additional strategy for infusing prior knowledge into tabular learning leverages the capabilities of large language models (LLMs)~\cite{anil2023palm,openai2023gpt4}. Trained over vast and heterogeneous collections of text, LLMs have shown impressive generalization, including on unfamiliar tasks and across disparate fields~\cite{brown2020language}. Recent investigations have focused on capitalizing on LLMs’ accumulated knowledge for tasks involving tabular data. This usually involves converting table-structured data into textual formats suitable for prompt-based LLM inputs~\cite{dinh2022lift,hegselmann2023tabllm,wang2023anypredict}. By embedding tabular content, descriptive features, and details regarding the predictive objective within prompts, models can reason over feature-task associations to generate predictions. Research directions in this area encompass specialized tuning procedures that modify only parts of the model such as LoRA~\cite{hu2021lora} and IA3~\cite{liu2022few}, as well as in-context learning, in which a limited number of annotated examples are introduced into the input prompt without any parameter updates during inference~\cite{dinh2022lift,hegselmann2023tabllm,nam2023semi,slack2023tablet}.

While the previous approaches have proven useful for enhancing few-shot tabular learning, the requirement to route every inference task through the LLM introduces practical obstacles for deployment in industrial settings. This research shifts away from the standard end-to-end, black-box LLM prediction pipeline, focusing instead on guiding the LLM to explicitly deduce the class-specific rules or conditions. With \textsf{FeatLLM}, these inferred rules are transformed into program code to derive fresh features, thus enabling predictions that bypass the LLM entirely. The model harnesses in-context learning to distill rules by leveraging both prior knowledge and the information contained within few-shot examples.

\section{Methods}
\textsf{FeatLLM} operates by extracting “rules”—the class-related conditions—from limited input samples, as illustrated in Figure~\ref{fig:main_model}. Using an LLM, these rules are parsed and converted into new binary feature representations for each sample, denoting the satisfaction or violation of each rule. These binary features serve as input to a dot product calculation with non-negative, trainable weights, yielding a prediction for class likelihood. Training this framework allows us to discern how each feature contributes to class probability (see Section~\ref{Sec:training}). To ensure robustness, the process employs repeated bagging and aggregates results through ensembling. The following sections detail the problem formulation and elaborate on the steps involved at each stage.

\vspace*{-0.12in}
\paragraph{Problem formulation.} Assume a tabular dataset comprised of $N$ labeled entries $\mathcal{D} = \{(\mathbf{x}^i, \mathbf{y}^i)\}_{i=1}^{N}$. Here, $\mathcal{D}$ contains $d$ features, meaning each $\mathbf{x}^i$ is a vector in $d$ dimensions. Feature names are expressed in natural language as $F = \{f_j\}_{j=1}^{d}$, with their corresponding definitions provided separately. Within the classification paradigm, each label $\mathbf{y}^i$ belongs to a pre-established set of classes. For $k$-shot learning evaluations, just $k$ ($<N$) randomly chosen labeled samples are used for training purposes.

\subsection{Prompt Design for Extracting Rules}
\label{Sec:prompt_design}
To facilitate rule extraction from the LLM via an expert-informed reasoning approach, we structure the prompting to emulate the methods employed by human specialists in tabular analysis tasks. The prompt is constructed with three principal parts (illustrated in Fig.~\ref{fig:prompt_for_rule} and Appendix~\ref{sec:example_rule}):

\vspace*{-0.12in}
\paragraph{Basic information description.} This segment supplies the key details needed for problem resolution (refer to the orange-highlighted text in Fig.~\ref{fig:prompt_for_rule}). Included are the task’s specification with associated label data, explanations of each feature, and examples for illustration. The task is introduced as a question format (e.g., ``Does this patient's myocardial infarction complications data indicate chronic heart failure? Yes or no?''), with more question templates available in Appendix~\ref{sec:task_desc}. Feature explanations specify data type and, when relevant, provide definitions; categorical features can be accompanied by sample categories. As illustrative demonstrations, a limited number of training samples are rendered into textual form with their true labels, following the serialization schema found in prior works~\cite{dinh2022lift,hegselmann2023tabllm}:

\begin{align}
&\text{Serialize}(\mathbf{x}^i,\mathbf{y}^i, F) = \nonumber \\ 
&``\ f_1\ \text{is}\ \mathbf{x}^i_1.\ ... \ f_d\  \text{is}\  \mathbf{x}^i_d.\ \ \text{Answer:}\  \mathbf{y}^i\ ”
\end{align}

\vspace*{-0.12in}
\paragraph{Reasoning instruction.} Our objective is to improve the LLM's capacity for reasoning by supplying targeted reasoning instructions (illustrated as blue text in Fig.~\ref{fig:prompt_for_rule}). These instructions start with an opening sentence aligned with the chain-of-thought paradigm~\cite{wei2022chain}. We then break down the inference process into two distinct phases. First, the LLM is tasked with determining causal associations or tendencies between the features and the task, drawing solely from its general knowledge or common sense and not from any provided demonstrations. This enables the model to utilize its pre-existing understanding of the problem domain. Second, the LLM integrates both exemplar demonstrations and insights gained in the initial phase to generate rules corresponding to each class. This structured reasoning sequence helps the model avoid false correlations from irrelevant columns and directs its attention toward more salient features. To mitigate risks such as extracting too few rules (which may cause underfitting) or generating an excessively lengthy prompt, we constrain the number of rules per class to a constant (specifically, 10)\footnote{An exploration of rule quantity effects can be found in Section~\ref{sec:analysis}.}. Furthermore, when composing rules, we instruct the LLM to reference the feature type details presented in the fundamental information segment, thereby reducing the likelihood of type inconsistencies within the generated rules.

\vspace*{-0.12in}
\paragraph{Response instruction.} To enable more efficient parsing and utilization of the derived rules, we explicitly guide the LLM in formatting its outputs with tailored response instructions (represented by the yellow text in Fig.~\ref{fig:prompt_for_rule}). The prompt delineates expectations for output format specific to each answer class (i.e., the response format) as well as the required structure for each individual rule (i.e., rule format). This guidance empowers the LLM to choose suitable formats depending on the feature type. Absent further instruction, the LLM has the flexibility to combine various rules through logical connectives such as AND or OR. A sample output is included in Appendix~\ref{sec:outcome-rule}.

\subsection{Inferring Class Likelihood via Rules}
\label{Sec:training}

\paragraph{Rule parsing for feature creation.} In this step, we convert the rules previously derived into new binary indicators. Each feature represents, for a given class, whether a data point adheres to the rule corresponding to that class, marked as '0' or '1'. These binary features serve as signals to infer the likelihood of class membership for each instance. Since the rules formulated by the LLM are expressed in natural language, an automated parsing strategy must be established to enable systematic feature extraction. Although we restrict the structure and phrasing of outputs and rules during their generation to simplify parsing, the LLM may still introduce minor syntactic variations~\cite{kovavc2023large}. For example, it might replace symbols such as ``$>$” and ``$<$” with their word equivalents, like ``greater than" and ``less than," introducing additional parsing complexity.

To mitigate parsing difficulties arising from textual noise, instead of employing intricate programming solutions, we make use of the LLM's own capabilities. The LLM excels at interpreting semantics even with syntactic discrepancies, and it can be prompted to generate succinct code solutions due to its training on coding data. Accordingly, our prompt includes the function signature, descriptions for inputs and outputs, and the relevant rules, which are supplied directly to the LLM. Figures~\ref{fig:prompt_for_parsing} and ~\ref{sec:example_parsing}-\ref{sec:example_parse_outcome} in the Appendix illustrate example prompts and their respective parsing outputs. The resulting code is then executed with Python’s \texttt{exec()} method to effectuate the data transformation.

\vspace*{-0.12in}
\paragraph{Class probability estimation.} With binary features established from rule evaluation for each class, one straightforward way to assess the likelihood that a data point belongs to a class is to tally the number of its rules met (i.e., summing the binary indicators per class). Nevertheless, each rule does not contribute equally, so their relative importance should be empirically determined using labeled training data. We address this by fitting a standard machine learning (ML) algorithm to the set of binary features per class, thus quantifying their importance. As a concrete example, consider applying a bias-free linear model. Define $\mathbf{z}_k^i \in \{0, 1\}^R$ as the vector of binary features generated from input $\mathbf{x}^i$ for class $k$, where $R$ represents the rule count per class, and $c$ denotes the total number of classes. The probability vector $\mathbf{p}^i$ for each class is given by:

\begin{align}
\text{logit}_k^i &= \max(\mathbf{w}_k, 0) \cdot \mathbf{z}_k^i, \nonumber \\
\mathbf{p}^i &= \text{Softmax}({[\text{logit}_{1}^i, ..., \text{logit}_{c}^i]}). \label{eq:infer_prob}
\end{align}

In this context, $\mathbf{w}_k$ are the adjustable weights within the linear framework, reflecting the significance assigned to each individual feature. By mapping these weights into the positive region, we guarantee that features crafted by the LLM cannot decrease the probability assigned to any class during the training phase. The logit vector undergoes transformation into class probabilities via the softmax operator. The learning process entails optimizing a cross-entropy objective to fit $\mathbf{w}_k$ using a limited set of training examples.

\vspace*{-0.12in}
\paragraph{Ensembling with bagging.} The entire pipeline is iterated multiple times to construct a set of inference models, whose outputs are aggregated through ensembling for final prediction. With the ensemble comprising $T$ independently learned weight vectors $\{\mathbf{w}[t]\}_{t=1}^T$, predictions are obtained by taking the mean of class probabilities $\mathbf{p}^i$ computed with each $\mathbf{w}[t]$, as described by Eq.~\ref{eq:infer_prob}.

Randomness in the rule generation procedure is induced by selecting a sufficiently high temperature setting during LLM inference. Additionally, we shuffle the sequence of few-shot examples presented to the LLM, considering findings that demonstrate how sequence ordering can alter LLM response~\cite{lu2022fantastically}\footnote{Further discussion regarding rule diversity is presented in Appendix~\ref{sec:diversity}}. To harness prediction diversity for improved robustness, we employ bagging, in which subsets of features or samples are drawn for each iteration—this approach becomes increasingly valuable as both feature and training sample sizes grow. The suggested ensemble methodology confers two principal benefits. Firstly, in instances where faulty rules emerge from certain trials, other trials producing correct rules help counterbalance them, thereby reinforcing system robustness. This concept resonates with the self-consistency phenomenon observed in LLMs~\cite{wang2022self}, suggesting that repeatedly inferred rules across different runs are more trustworthy. Secondly, ensembling alleviates the constraint imposed by the LLM’s input prompt length: when tabular data surpass the input limit, bagging assists in managing this issue, sustaining reliable performance. Though an individual set of rules may fail to comprehensively capture all feature interrelations, the collective output from multiple weak learners—each generated from separate trials—mitigates deficits in single-trial coverage of features or instances.

\section{Experiments}
We assess the effectiveness of \textsf{FeatLLM} across a selection of tabular datasets and undertake a series of analyses to explore the operational characteristics of our framework. Owing to limitations of space, supplementary results—including investigations using alternative LLMs, qualitative assessments, and other auxiliary experiments—are provided in the Appendix.

\subsection{Performance Evaluation}
\paragraph{Datasets.}
Our empirical study spans 13 datasets geared towards binary or multi-class classification, namely: (1) Adult~\cite{asuncion2007uci}, which identifies individuals earning above \$50,000 each year; (2) Bank~\cite{moro2014data}, forecasting customers likely to opt for a term deposit; (3) Blood~\cite{yeh2009knowledge}, detecting donors expected to return for future donations; (4) Car~\cite{kadra2021well}, rating the quality of vehicles; (5) Communities~\cite{misc_communities_and_crime_183}, estimating the crime rates for specified regions; (6) Credit-g~\cite{kadra2021well}, determining if a subject represents favorable or unfavorable credit risk; (7) Diabetes\footnote{\texttt{kaggle.com/datasets/uciml/pima-indians-diabetes-database}}, which predicts diabetes onset for individuals; (8) Heart\footnote{\texttt{kaggle.com/datasets/fedesoriano/heart-failure-prediction}}, projecting risk for coronary artery disease; and (9) Myocardial~\cite{misc_myocardial_infarction_complications_579}, forecasting chronic heart failure occurrences.

To further diversify the evaluation, we include newly released and synthetic datasets that remain largely absent from prior LLM training: (10) Cultivars, which infers expected grain yield for soybean cultivars\footnote{\texttt{archive.ics.uci.edu/dataset/913}}; (11) NHANES, used to classify individuals' age groups based on their records\footnote{\texttt{archive.ics.uci.edu/dataset/887}}; (12) Sequence-type, categorizing sequences into four classes—arithmetic, geometric, Fibonacci, or Collatz; and (13) Solution-mix, which determines if a mixture's concentration from four solutions surpasses 0.5. The Cultivars and NHANES datasets, added to the UCI Machine Learning Repository post-September 2023, were unavailable during pre-training of the LLM API (gpt-3.5-turbo-0613), eliminating memorization as a concern for our model. The Sequence-type and Solution-mix are synthetic constructs, ensuring exclusion from the LLM training corpus and thus are suitable for unbiased evaluation.

A detailed summary of dataset scale and complexity is found in Table~\ref{Tab:basic-desc}.
Each dataset features explicit attribute naming and explanation. Some, such as `Communities' and `Myocardial', encompass in excess of 100 attributes, challenging context capacity constraints inherent in LLMs.

\vspace*{-0.12in}
\paragraph{Baselines.}
We benchmark against nine distinct baseline methods. The initial three comprise conventional approaches to tabular data: (1) Logistic regression (LogReg), (2) XGBoost~\cite{chen2016xgboost}, and (3) RandomForest~\cite{ho1995random}. The following two leverage unlabeled data resources: (4) SCARF~\cite{bahri2022scarf}, and (5) STUNT~\cite{nam2023stunt}, although the assumption of access to large volumes of unlabeled data rarely holds in practical settings. We further include (6) TabPFN~\cite{hollmann2022tabpfn}, which pretrains a model via a diverse, synthetically generated dataset. The remaining approaches rely on LLM-based methodologies: (7) In-context learning (In-context)~\cite{wei2022emergent}, which integrates few-shot examples directly in the prompt with no additional model training; (8) TABLET~\cite{slack2023tablet}, enriching the prompt by appending external classifier-generated hints, such as rule sets and prototypes, to boost inference reliability; and (9) TabLLM~\cite{hegselmann2023tabllm}, which uses parameter-efficient fine-tuning approaches such as IA3~\cite{liu2022few} to adapt the LLM with limited few-shot data.

\vspace*{-0.12in}
\paragraph{Implementation details.}
\textsf{FeatLLM} adopts GPT-3.5 as its LLM backbone, although the framework maintains flexibility to incorporate various LLMs (refer to Appendix~\ref{sec:backbones} for comparative results across backbones). For LLM inference, the temperature parameter is fixed at 0.5, while top-p remains set to the API default of 1. Both the ensemble size and the count of extraction rules are set to 20 and 10, respectively. Further analysis regarding hyper-parameter selection can be found in Figure~\ref{fig:hyperparameter2} and Appendix~\ref{sec:hyper-parameter}. 
The linear model leverages the Adam optimizer with a learning rate of 0.01 and is trained over 200 epochs. To determine the optimal stopping point, $k$-fold cross-validation is applied. For baseline reproduction, GPT-3.5 is utilized for in-context learning; for TabLLM, T0~\cite{sanh2022multitask} is chosen as specified in the original implementation. It is important to note that TabLLM imposes a limitation in that it only allows LLMs with available public checkpoints, making GPT-3.5 unavailable for use within this baseline. Parameter selection for traditional models (LogReg, XGBoost, RandomForest) is accomplished through grid-search coupled with $k$-fold cross-validation. The value of $k$ is set as either 2 or 4, ensuring every class is present in the training split. For supplementary implementation specifics, see Appendix~\ref{sec:baselines}.

\vspace*{-0.12in}
\paragraph{Results.}
Tables~\ref{tab:main1} and \ref{tab:main2} summarize performance outcomes for a diverse set of datasets. Each experiment is conducted three times with distinct random splits between training and testing sets, with mean and standard deviation of the resulting AUC values reported. \textsf{FeatLLM} persistently outperforms other baseline methods or attains the second-highest scores. Remarkably, our framework delivers performance comparable to standard tabular baselines that utilize 32 shots, even when operating with only 4 shots (refer to Fig.~\ref{fig:compares}). While STUNT and TabPFN exhibit pronounced performance gains in specific datasets, their overall superiority to conventional machine learning baselines remains limited. Additionally, LLM-powered approaches—including in-context learning, TABLET, and TabLLM—were unable to handle inference tasks on tabular datasets characterized by a large number of features. By contrast, our method, leveraging both bagging and ensemble strategies, maintained robust performance even as feature dimensionality increased. Notably, the principles of bagging and ensembling that underpin our framework could, in theory, be incorporated into other LLM-centric methods; however, such an extension would require a twofold increase in inference computations for each instance, greatly exacerbating computational demands and rendering the approach infeasible.

\subsection{Analysis \& Discussion}
\label{sec:analysis}
\paragraph{Ablation study.} 
We investigate the roles of key components via an ablation analysis targeting the following: (1) \textsf{-Tuning}: bypassing the weight adjustment stage within the linear model and directly summing the binary feature outputs to compute logits; (2) \textsf{-Ensemble}: removing the ensemble aggregation process; (3) \textsf{-Description}: excluding feature-level textual descriptions; and (4) \textsf{-Reasoning}: eliminating the Step-1 reasoning instructions during rule generation.

Table~\ref{tab:ablation} reports the average changes in AUC across all datasets due to these ablations. Modifying or removing any single component consistently led to diminished performance. Notably, the impact of the tuning module escalates with increased shot count, implying that accurate estimation of rule significance is particularly advantageous when ample labeled data are available. Conversely, the absence of feature descriptions and reasoning cues causes greater drops in performance when few shots are present, underlining the importance of LLMs’ prior knowledge in low-data regimes. Finally, the ensemble approach proposed herein contributes consistent and robust gains under all experimental configurations.

\vspace*{-0.12in}
\paragraph{Training \& inference time comparison.} We evaluate and contrast the training and inference durations across several approaches. All experiments are carried out on the Adult dataset, utilizing 16 training samples ("shots"). An A100 GPU is standard for these tests, except for TabLLM, which employs four A100 GPUs to facilitate model parallelism. Table~\ref{tab:cost} details the timing results. A key observation is that \textsf{FeatLLM} maintains inference times nearly equivalent to traditional methods. By comparison, inference times for other LLM-driven techniques—In-context learning, TABLET, and TabLLM—are substantially longer. Regarding training, \textsf{FeatLLM}'s time requirement aligns closely with other LLM-based strategies: it necessitates only 30 API calls, which do not represent a notable financial or computational overhead and are amenable to parallel execution.

\vspace*{-0.12in}
\paragraph{Quality of features generated by \textsf{FeatLLM}.}
To assess how informative the rule-based features from \textsf{FeatLLM} are for real-world tasks, we perform a simple experiment. Specifically, we calculate the average AUROC value between each newly produced feature and the target variable, then determine the proportion of features with AUROC exceeding 0.5 for each dataset. Table~\ref{tab:quality} summarizes these findings, demonstrating that most generated rules substantially relate to the target attribute and contribute meaningful information toward solving the target task (see the “Before tuning” column). Additionally, as the linear model is tuned with data, rules that lack significant relevance—those associated with weights below a certain threshold—are systematically removed (see “After tuning” column). Here, a threshold value of 0.1 is chosen, guided by the weight distribution histogram.

\vspace*{-0.12in}
\paragraph{Effect of adding spurious correlations.} We investigated how well different approaches eliminate spurious, irrelevant correlations when working with limited labeled data. For this purpose, we performed an experiment using the Heart dataset, augmenting it by randomly appending columns from the Adult dataset that have no connection to the prediction objective of the Heart dataset. To observe the impact, we tracked how each method’s accuracy shifted as more unrelated columns were added. Ensuring a consistent baseline, all models had access to 8 labeled samples and no unlabeled data; consequently, approaches such as SCARF and STUNT, which rely on unlabeled samples, were not included in our evaluation.

Figure~\ref{fig:spurious} presents how performance changes as spurious correlations increase. Among the methods, \textsf{FeatLLM} exhibited minimal decline in accuracy, outperforming other baselines under the influence of noisy features. This resilience is likely due to the explicit integration of domain knowledge in the rule extraction process, which compels \textsf{FeatLLM} to focus on the relevance between task-specific objectives and input features—effectively curbing rules that depend on random or irrelevant columns, thereby maintaining model robustness. Supporting this, the proportion of rules built upon noise features was observed to be limited to only 1-2\% of the total. However, we also noted that when additional columns were derived from datasets closely aligned with the task, \textsf{FeatLLM} occasionally picked up and overfit to misleading correlations, mistakenly conflating features with actual causality (see Appendix~\ref{sec:spurious-appendix}).

\vspace*{-0.12in}
\paragraph{Hyper-parameter analysis}
There are two primary hyper-parameters associated with \textsf{FeatLLM}: the ensemble size, and the quantity of rules per class to extract from each LLM prompt. To assess the sensitivity of model performance to these settings, we performed empirical evaluations. Findings reveal that increasing the ensemble count enhances performance, but after approximately 20 ensembles, the improvements plateau (refer to Fig.~\ref{fig:appendix-hyperparameter1} in the Appendix). Regarding rule count per class, our results indicate no statistically significant effect; however, selecting 10 rules achieves the best balance between prompt complexity and predictive accuracy (see Fig.~\ref{fig:hyperparameter2}). We hypothesize that expanding rule numbers boosts input dimensionality, infusing more information but also introducing overfitting risks, particularly in low-shot contexts.

\section{Conclusion}
We introduce \textsf{FeatLLM}, a method that sets a new benchmark in leveraging LLMs for few-shot tabular tasks. Rather than directly predicting outcomes with LLMs, \textsf{FeatLLM} harnesses these models to derive predictive criteria akin to feature engineering. The process includes rule extraction and bagged ensemble methods, resulting in negligible inference burden, compatibility with API usage alone (no model retraining), and effective handling regardless of feature space scale. In practice, \textsf{FeatLLM} delivers robust accuracy, making it a compelling, data-efficient alternative for real-world tabular datasets.

\textsf{FeatLLM} is specifically optimized for settings with scarce training data. Looking ahead, we intend to enhance the framework’s capacity for automated feature generation in contexts with more abundant data, incorporate more diverse forms of features beyond rule-based, and prioritize interpretability to broaden its applicability across various use cases.

\section{Broader Impact}
The \textsf{FeatLLM} framework seeks to exploit the substantial prior knowledge inherent in LLMs to excel at tabular learning tasks—surpassing conventional supervised methods, especially under limited data conditions. By infusing pretrained knowledge, the approach meaningfully addresses overfitting in low-data regimes. This is particularly advantageous for practitioners in sectors such as finance and healthcare, where labeled data acquisition is costly or logistically difficult and pretrained LLMs may already encode domain-relevant expertise. As the framework develops towards wider implementation, a critical future research direction will be investigating how societal biases or misinformation present in LLM prior knowledge might influence or override model predictions—even relative to empirical data—thus ensuring responsible adoption.

\end{document}